\chapter{Benefits, Conclusion and Future Work}

\section{Benefits and Real-World Applications}
\begin{table}[H]
\centering
\small
\caption{Supply-chain impact of SmartDock}
\label{tab:benefits}
\begin{tabularx}{\textwidth}{>{\bfseries\raggedright\arraybackslash}p{2.6cm}X}
\toprule
Dimension & Impact \\
\midrule
Cost & Less truck waiting means lower detention charges and driver overtime; labour can be rostered to the known appointment load (heatmap, busiest hours) instead of covering unpredictable peaks; free and open-source software avoids licence fees. \\
Visibility & Both sides see the same schedule in real time: carriers see free capacity and status; managers see every truck's position (expected, in yard, at door, departed) and live KPIs. \\
Speed & Self-service booking replaces phone and e-mail; the simulation indicates that yard waiting falls from about 19 minutes to about 1 minute per truck; waitlist auto-assignment refills cancelled capacity without manual effort. \\
Traceability & Every appointment carries PO/ASN references, timestamps for each yard event and a complete audit trail, supporting detention disputes, compliance and root-cause analysis. \\
Reliability & Conflict-free scheduling (including dock compatibility and hazmat safety) and a carrier reliability score that creates accountability and supports carrier selection. \\
Sustainability & Shorter queues mean less engine idling and local emissions near the site \cite{giuliano2007}. \\
\bottomrule
\end{tabularx}
\end{table}

\paragraph{Applications.} Distribution centres of retailers and e-commerce companies (inbound supplier deliveries, outbound store or courier pickups); cold-chain warehouses (reefer-only doors and pharma 2--8\,$^{\circ}$C loads); chemical and fuel depots (hazmat bays); third-party logistics (3PL) providers running several client warehouses on one platform (multi-tenancy); manufacturing plants receiving just-in-time inbound parts; and, with minor changes, container freight stations and inland container depots.

\section{Conclusion}
This project designed, implemented, tested and deployed \textbf{SmartDock}, a multi-tenant dock appointment scheduling and yard management system. Mapping back to the objectives:
\begin{itemize}
  \item Multi-tenancy with per-warehouse carrier approval and strict tenant isolation was implemented and verified by automated tests.
  \item A configurable scheduling engine supports slot lengths of 15--90 minutes, turnover buffers, holidays and maintenance windows across time zones.
  \item Loads are matched to doors by direction, trailer type and hazmat certification.
  \item Slot capacity is guaranteed under concurrency across multiple servers; a controlled experiment confirmed the guarantee and demonstrated the failure of the naive approach.
  \item Recurring series, rescheduling and an automatic FIFO waitlist cover realistic carrier needs.
  \item Yard events are digitised, late arrivals and no-shows are detected automatically, and carriers receive a reliability score.
  \item Managers obtain KPIs, heatmaps, scorecards, CSV/PDF reports and a gate sheet; every action is audited and notified.
  \item The system is live on the cloud and passes 58 end-to-end checks.
\end{itemize}
SmartDock shows that a focused information system at the warehouse--transport interface can convert an uncontrolled, phone-based process into a measured, collaborative one, improving cost, visibility, speed and traceability in the supply chain.

\section{Project Links}
\begin{tabular}{@{}ll@{}}
\textbf{Live system} & \url{\LiveURL} \\
\textbf{Source code (GitHub)} & \url{\GitHubURL} \\
\textbf{Demo logins} & \texttt{admin@smartdock.demo}, \texttt{carrier@smartdock.demo} (password \texttt{Demo@123}) \\
\textbf{Run locally} & \texttt{npm run dev:demo} (server) after building the client \\
\textbf{Run tests} & \texttt{npm run verify:critical} (58 end-to-end checks) \\
\end{tabular}

\section{Future Enhancements}
\begin{itemize}
  \item \textbf{AI/ML:} predict no-show probability and arrival delay per carrier, lane and time of day; recommend slots and automatically overbook slightly where no-shows are likely; forecast daily inbound volume for labour planning.
  \item \textbf{Optimisation:} automatic assignment of trucks to doors (minimising yard wait and travel inside the building) and load-dependent slot lengths.
  \item \textbf{IoT and computer vision:} automatic gate check-in by number-plate recognition, RFID or GPS geofencing; dock-door sensors for real dock start and end times; reefer temperature monitoring.
  \item \textbf{Integration:} import purchase orders and ASNs from ERP/WMS, push appointments to TMS, and provide webhooks and a public API for carriers.
  \item \textbf{Mobile app:} a driver app with turn-by-turn directions to the assigned door, digital proof of delivery and live status; installable as a progressive web app.
  \item \textbf{Commercial features:} detention and no-show fee calculation, carrier contracts and service-level agreements, multi-language support.
  \item \textbf{Scalability and operations:} real-time updates with WebSockets instead of polling, caching of analytics, a super-administrator console for 3PL operators, horizontal scaling on a paid plan, and automated CI tests on every commit.
\end{itemize}
